%% FullCourse_10Lecs -- Lecture 7, Act 3 (ATTEND, part 1/2): The Attention Mechanism
%% ch14 rebuild (2026-07-07); worked-example-first reorder (2026-07-09).
%% Reorganized top-down per user request: open with a worked example + the goal,
%% then the approach (Q/K/V roles + retrieval intuition), then WHERE the
%% projection matrices come from (the training objective -- promoted from the
%% old end-of-deck "How Q/K/V Are Learned" frame), then the mechanics step by
%% step (which now carry the worked example's numbers, absorbing the old
%% standalone "Worked Example Part 1/2"). Frame bodies reuse the ch14-rebuild
%% content; the self-attention worked example's scores were recomputed so that
%% softmax(scores) reproduces the 0.50/0.40 headline used by the "Q/K/V at a
%% Glance" diagram (the old 2.3/1.9/1.1/... table softmaxed to ~0.45/0.30/0.14,
%% not the stated 0.50/0.40/0.05 -- now internally consistent).
%% Original ch14-rebuild frames were copied from
%%   archive_pre_split/pre_ch14_rebuild_2026-07-07/Lec07_T2a_Attention_Mechanism.tex
%% The Q/K/V-at-a-glance diagram keeps the corrected arrow semantics (Q-K match
%% sets the weight, the weight applies to V).

%% Shared preamble for the FullCourse_10Lecs topic decks.
%%
%% Each topic .tex \input's this file, then defines its own \title, \date
%% override (if any), \graphicspath, and \begin{document}.
%%
%% Reconstructed verbatim from the inlined copy preserved in
%% Lec10_Case_Studies/Lec10_T8_Case6_DeepRamsey.tex (lines 23-190), which is
%% explicitly annotated there as "from ../shared_preamble.tex, copied verbatim".
%% Base preamble matches MiniCourse_8hr/Slides/shared_preamble.tex; this version
%% additionally provides \sectiondivider, the conceptbox/cautionbox/examplebox/
%% takeawaybox tcolorboxes, and \compacttables used across the split decks.

\documentclass[10pt,english,aspectratio=169]{beamer}

\usepackage{ae,aecompl}
\usepackage[T1]{fontenc}
\usepackage{textcomp}
\usepackage[utf8]{inputenc}
\usepackage{babel}
\usepackage{datetime2}

\setcounter{secnumdepth}{3}
\setcounter{tocdepth}{3}

\usepackage{amsmath, amssymb, amsthm, graphicx}
\usepackage{booktabs, multirow}
\usepackage{tikz}
\usepackage{pgfplots}
\pgfplotsset{compat=1.16}
\usepackage[most]{tcolorbox}
\tcbuselibrary{skins, breakable, theorems}
\usepackage{listings}
\usepackage{xcolor}

\lstset{
  basicstyle=\ttfamily\small,
  keywordstyle=\color{blue!80!black}\bfseries,
  commentstyle=\color{green!50!black}\itshape,
  stringstyle=\color{orange!80!black},
  backgroundcolor=\color{gray!8},
  frame=single,
  rulecolor=\color{gray!40},
  breaklines=true,
  showstringspaces=false,
  numbers=none,
}

\lstdefinestyle{pythonstyle}{
    language=Python,
    basicstyle=\ttfamily\footnotesize,
    keywordstyle=\color{blue!80!black}\bfseries,
    commentstyle=\color{green!50!black}\itshape,
    stringstyle=\color{orange!80!black},
    frame=single,
    breaklines=true,
    showstringspaces=false,
    numbers=left,
    numbersep=5pt,
    numberstyle=\tiny\color{gray},
    backgroundcolor=\color{gray!8},
    rulecolor=\color{gray!40}
}

\ifx\hypersetup\undefined
  \AtBeginDocument{%
    \hypersetup{bookmarks=false,breaklinks=true,pdfborder={0 0 0},colorlinks=true,
      linkcolor=DarkText,urlcolor=RedTitle}%
  }
\else
  \hypersetup{bookmarks=false,breaklinks=true,pdfborder={0 0 0},colorlinks=true,
    linkcolor=DarkText,urlcolor=RedTitle}%
\fi

\makeatletter
\newcommand\makebeamertitle{%
  \begin{frame}[plain]
    \vspace*{1.0cm}
    \centering
    {\usebeamerfont{title}\usebeamercolor[fg]{title}\inserttitle\par}
    \vspace{1.0cm}
    {\usebeamerfont{author}\usebeamercolor[fg]{author}\insertauthor\par}
    \vspace{0.2cm}
    {\usebeamerfont{date}\usebeamercolor[fg]{date}\insertdate\par}
  \end{frame}}%
\AtBeginDocument{%
  \let\origtableofcontents=\tableofcontents
  \def\tableofcontents{\@ifnextchar[{\origtableofcontents}{\gobbletableofcontents}}
  \def\gobbletableofcontents#1{\origtableofcontents}%
}

\usetheme{metropolis}
\usefonttheme{professionalfonts}

\definecolor{LightGray}{RGB}{230,230,230}
\definecolor{RedTitle}{RGB}{0,0,200}
\definecolor{VeryLightGray}{RGB}{245,245,245}
\definecolor{DarkText}{RGB}{0,0,0}

\setbeamercolor{background canvas}{bg=white}
\setbeamercolor{title}{fg=RedTitle}
\setbeamercolor{author}{fg=DarkText}
\setbeamercolor{date}{fg=DarkText}
\setbeamercolor{frametitle}{bg=VeryLightGray,fg=DarkText}
\setbeamercolor{normal text}{fg=DarkText}
\setbeamerfont{title}{size=\large}

\setbeamersize{text margin left=5mm, text margin right=5mm}
\addtobeamertemplate{frametitle}{}{\vspace{-0.5em}}
\newcommand{\reducevspace}{\vspace{-0.5cm}}

% Shared section divider and box styles for split topic decks.
\newcommand{\sectiondivider}[2]{%
\begin{frame}[plain,noframenumbering]
\begin{center}
\vfill
{\Large\textbf{#1}\par}
\vspace{0.35cm}
{\normalsize #2\par}
\vfill
\end{center}
\end{frame}
}

\newtcolorbox{conceptbox}[1][]{
  colback=blue!3,
  colframe=RedTitle!55,
  boxrule=0.35mm,
  arc=1mm,
  left=2mm,
  right=2mm,
  top=1mm,
  bottom=1mm,
  #1
}
\newtcolorbox{cautionbox}[1][]{
  colback=red!3,
  colframe=red!50!black,
  boxrule=0.35mm,
  arc=1mm,
  left=2mm,
  right=2mm,
  top=1mm,
  bottom=1mm,
  #1
}
\newtcolorbox{examplebox}[1][]{
  colback=green!3,
  colframe=green!45!black,
  boxrule=0.35mm,
  arc=1mm,
  left=2mm,
  right=2mm,
  top=1mm,
  bottom=1mm,
  #1
}
\newtcolorbox{takeawaybox}[1][]{
  colback=VeryLightGray,
  colframe=RedTitle!55,
  boxrule=0.35mm,
  arc=1mm,
  left=2mm,
  right=2mm,
  top=1mm,
  bottom=1mm,
  #1
}
\newcommand{\compacttables}{\renewcommand{\arraystretch}{1.15}}

\setbeamertemplate{navigation symbols}{}
\setbeamertemplate{footline}{%
  \hfill%
  \usebeamercolor[fg]{page number in head/foot}%
  \usebeamerfont{page number in head/foot}%
  \insertframenumber\,/\,\inserttotalframenumber\kern1em\vskip2pt%
}

\makeatother

\author{Zhigang Feng}
% "date with time": compile timestamp so each build is identifiable.
\date{\today\ \DTMcurrenttime}


\metroset{sectionpage=none}

\graphicspath{{./pic/}{../../MiniCourse_8hr/Slides/Module3_LLM_Text/pic/}{../}}

\usetikzlibrary{arrows.meta,positioning,shapes,calc,fit,backgrounds}

%% Lec07 shared MAP slide: "one map, five acts" recurring diagram.
%% Adapted from ../Lec09_Agentic_AI/lec09_map.tex (\LecNineMap -> \LecSevenMap).
%% Input this file AFTER ../shared_preamble.tex (needs RedTitle, VeryLightGray)
%% and call \LecSevenMap{k} inside a frame, where k in {1..5} highlights the
%% current act ("you are here") and k = 0 renders the closing reprise with all
%% five acts complete (used by T5b's closing frame).
%% Filename deliberately does not match the build glob Lec*_T*.tex, so the
%% build script never tries to compile it standalone.

\newcommand{\LecSevenMap}[1]{%
% Precompute per-box style selectors; TikZ option lists are comma-split before
% TeX conditionals run, so \ifnum must not sit inside a [...] options list.
\ifnum#1=0
  \tikzset{lecmapa/.style={lecmapdone}, lecmapb/.style={lecmapdone},
           lecmapc/.style={lecmapdone}, lecmapd/.style={lecmapdone},
           lecmape/.style={lecmapdone}}%
\else
  \tikzset{lecmapa/.style={}, lecmapb/.style={}, lecmapc/.style={},
           lecmapd/.style={}, lecmape/.style={}}%
  \ifnum#1=1 \tikzset{lecmapa/.style={lecmapcur}}\fi
  \ifnum#1=2 \tikzset{lecmapb/.style={lecmapcur}}\fi
  \ifnum#1=3 \tikzset{lecmapc/.style={lecmapcur}}\fi
  \ifnum#1=4 \tikzset{lecmapd/.style={lecmapcur}}\fi
  \ifnum#1=5 \tikzset{lecmape/.style={lecmapcur}}\fi
\fi
\begin{center}
\begin{tikzpicture}[
    act/.style={rectangle, rounded corners, draw=black!60, fill=VeryLightGray,
        minimum width=2.6cm, minimum height=0.92cm, text width=2.5cm,
        align=center, font=\scriptsize, text=DarkText},
    lecmapcur/.style={draw=RedTitle, very thick, fill=orange!20},
    lecmapdone/.style={draw=green!45!black, thick, fill=green!10},
    q/.style={font=\tiny\itshape, text width=2.7cm, align=center, text=black!70},
    sat/.style={rectangle, rounded corners, draw=black!45, dashed, fill=white,
        text width=5.0cm, align=center, font=\tiny, text=black!70,
        minimum height=0.5cm},
    barlab/.style={font=\tiny\bfseries, text=black!60},
    arr/.style={-{Stealth[length=2mm]}, thick, gray!60}
]
    % --- five act boxes ---
    \node[act, lecmapa] (a1) at (0,0)
        {\textbf{7.1 FRAME}\\ \tiny text as measurement};
    \node[act, lecmapb] (a2) at (2.85,0)
        {\textbf{7.2 REPRESENT}\\ \tiny tokens $\to$ vectors};
    \node[act, lecmapc] (a3) at (5.7,0)
        {\textbf{7.3 ATTEND}\\ \tiny attention \& Transformer};
    \node[act, lecmapd] (a4) at (8.55,0)
        {\textbf{7.4 PRETRAIN}\\ \tiny LMs \& alignment};
    \node[act, lecmape] (a5) at (11.4,0)
        {\textbf{7.5 MEASURE}\\ \tiny limits, apps, validation};
    \draw[arr] (a1) -- (a2);
    \draw[arr] (a2) -- (a3);
    \draw[arr] (a3) -- (a4);
    \draw[arr] (a4) -- (a5);
    % --- the student question each act answers ---
    \node[q] at (0,-0.82)    {what generates\\ this text?};
    \node[q] at (2.85,-0.82) {how do we turn it\\ into numbers?};
    \node[q] at (5.7,-0.82)  {how does context\\ get encoded?};
    \node[q] at (8.55,-0.82) {where does a pretrained\\ model come from?};
    \node[q] at (11.4,-0.82) {is my measure valid ---\\ and defensible?};
    % --- "you are here" marker above the current act ---
    \ifnum#1=1 \node[font=\scriptsize\bfseries, text=RedTitle] at (0,0.95)
        {you are here\ $\blacktriangledown$};\fi
    \ifnum#1=2 \node[font=\scriptsize\bfseries, text=RedTitle] at (2.85,0.95)
        {you are here\ $\blacktriangledown$};\fi
    \ifnum#1=3 \node[font=\scriptsize\bfseries, text=RedTitle] at (5.7,0.95)
        {you are here\ $\blacktriangledown$};\fi
    \ifnum#1=4 \node[font=\scriptsize\bfseries, text=RedTitle] at (8.55,0.95)
        {you are here\ $\blacktriangledown$};\fi
    \ifnum#1=5 \node[font=\scriptsize\bfseries, text=RedTitle] at (11.4,0.95)
        {you are here\ $\blacktriangledown$};\fi
    \ifnum#1=0 \node[font=\scriptsize\bfseries, text=green!45!black] at (5.7,0.95)
        {all five acts complete};\fi
    % --- bottom band: the measurement sandwich ---
    \fill[blue!10]   (-1.3,-1.55) rectangle (1.40,-1.22);
    \fill[green!10]  (1.65,-1.55) rectangle (9.85,-1.22);
    \fill[orange!15] (10.1,-1.55) rectangle (12.7,-1.22);
    \node[barlab] at (0.05,-1.385) {WHY TEXT = MEASUREMENT};
    \node[barlab] at (5.75,-1.385) {HOW THE MACHINE READS TEXT};
    \node[barlab] at (11.4,-1.385) {MEASURE, VALIDATE, DISCLOSE};
    % --- the recurring spine (text DGP), printed under the band ---
    \node[font=\tiny\itshape, text=black!60] at (5.7,-1.83)
        {the spine:\ \ $x_i = g(s_i,\,a_i,\,c_i)+\varepsilon_i
         \;\longrightarrow\; m(x_i)
         \;\longrightarrow\; y_i = \beta\,m(x_i)+\gamma z_i+u_i$};
    % --- dashed satellites: the neighbouring lectures ---
    \node[sat] at (2.0,-2.42)
        {\textit{before 7.1:} Lec03--04 gave you neural nets, backprop \& PyTorch --- Lec07 teaches them to read};
    \node[sat] at (9.8,-2.42)
        {\textit{after 7.5:} Lec08 RAG (grounded memory) $\to$ Lec09 agents (grounded action) $\to$ Lec10 cases};
\end{tikzpicture}
\end{center}
}


\title[AI for Econ Research]{AI for Economic Research: Dynamic Models, Language, and Agents\\
Lecture 7.3a: The Attention Mechanism}

\begin{document}

\makebeamertitle

%=============================================================================
\begin{frame}{Lecture 7 in One Map}
\LecSevenMap{3}
\medskip
\textit{\footnotesize Route: one worked example $\to$ the goal (a context vector per token) $\to$ the approach (Q/K/V) $\to$ where the matrices come from (the training objective) $\to$ the mechanics, step by step $\to$ multi-head; part 2/2 (T3b) wraps the head into the full Transformer layer.}
\end{frame}

%=============================================================================
\section{The Goal, Through One Example}
\sectiondivider{The Goal, Through One Example}{One ambiguous word --- and the context-aware vector we want for it}

\begin{frame}{Start Here: One Word, One Problem}
\small
\textbf{A sentence from a policy report:} ``The central bank raised interest rates amid rising inflation.''

\smallskip
\textbf{Tokenization} (9 word tokens, numbered left to right):
\[
\underbrace{\text{The}}_{1}\;\underbrace{\text{central}}_{2}\;\underbrace{\text{bank}}_{3}\;\underbrace{\text{raised}}_{4}\;\underbrace{\text{interest}}_{5}\;\underbrace{\text{rates}}_{6}\;\underbrace{\text{amid}}_{7}\;\underbrace{\text{rising}}_{8}\;\underbrace{\text{inflation}}_{9}
\]

\begin{itemize}
\item \textbf{The problem.} In isolation, \emph{interest} (position 5) is ambiguous --- monetary-policy \emph{rates}, personal \emph{curiosity}, or a financial \emph{stake}? We want a vector $e'^{(5)}$ encoding how it is used \emph{in this sentence}.
\item \textbf{What a reader does.} Seeing \emph{interest}, you ask ``interest in \emph{what}?'' --- your eye jumps to \emph{rates} and \emph{inflation}, the words that fix its meaning.
\item \textbf{What attention does (previewed).} Exactly that, numerically. Of all the attention \emph{interest} pays across the sentence:
\[
   \text{attention on \emph{rates}} \approx 0.50, \quad
   \text{on \emph{inflation}} \approx 0.40, \quad
   \text{on all else} \approx 0.10,
\]
then it \emph{blends their content} into $e'^{(5)}$ --- landing it squarely in monetary-policy territory.
\end{itemize}

\smallskip
\begin{takeawaybox}
\textbf{Those weights were never hand-coded --- no dictionary said ``\emph{interest} goes with \emph{rates}.''} How the model \emph{finds} them, and what $0.50/0.40$ computes, is the rest of this deck.
\end{takeawaybox}
\end{frame}

\begin{frame}{What We're Trying to Achieve}
\begin{itemize}
\item Zoom out from one word. A Transformer is a function
\[
   \mathcal{T}\colon \bigl(\mathbb{R}^{d}\bigr)^{n} \;\longrightarrow\; \bigl(\mathbb{R}^{d}\bigr)^{n},
   \quad \bigl[e^{(1)}, \ldots, e^{(n)}\bigr] \mapsto \bigl[e'^{(1)}, \ldots, e'^{(n)}\bigr],
\]
mapping $n$ input vectors to $n$ contextualized outputs of the same dimension. (Inputs are token embeddings $+$ positional encodings; details in T3b.)
\smallskip
\item \textbf{The single conceptual jump: static $\to$ contextual.} Word2Vec gave \emph{one vector per word type} --- \emph{bank} in ``river bank'' and ``central bank'' share it. A Transformer gives \emph{one vector per token occurrence}: the two \emph{bank}s separate. Our $e'^{(5)}$ for \emph{interest} is exactly such a vector.
\smallskip
\item \textbf{What ``contextualized'' packs in.} Each $e'^{(i)}$ encodes at once (a) the identity of $t^{(i)}$, (b) its role in \emph{this} sequence, (c) a weighted blend of every other token's contribution.
\smallskip
\item \textbf{Not the final answer.} These are an \emph{intermediate representation}. A small \emph{task head} sits on top: project $e'^{(n)}$ onto the vocabulary (generation) or $e'^{(\text{[CLS]})}$ onto labels (classification).
\end{itemize}
\end{frame}

%=============================================================================
\section{The Approach: Query, Key, Value}
\sectiondivider{The Approach: Query, Key, Value}{Score with queries and keys, read with values --- a soft database lookup}

\begin{frame}{The Approach: Score, Normalize, Blend}
\begin{itemize}
\item For each position $i$, ask: \emph{``which other positions are most relevant to $i$?''} Then build $i$'s new representation as a \emph{weighted sum} over all positions:
\[
   \tilde e^{(i)} \;=\; \sum_{j} \alpha_{j}^{(i)}\,(\text{content of } j), \qquad
   \alpha_{j}^{(i)} \ge 0, \quad \sum_{j}\alpha_{j}^{(i)} = 1.
\]
The weights $\alpha^{(5)} = (\ldots, 0.50, \ldots, 0.40, \ldots)$ from the previous slide are exactly these.

\medskip
\item \textbf{Crucial properties:}
\begin{itemize}
    \item Every token attends to every other in \emph{one} step --- no recurrence.
    \item Weights are computed \emph{dynamically} from the input, not fixed.
    \item All positions are processed in \emph{parallel} $\Rightarrow$ trains efficiently on GPUs.
\end{itemize}

\medskip
\item Repeat this ``everyone looks at everyone'' operation $L$ times (e.g.\ $L=96$ in GPT-3) with different learned weights, and the model builds deeply contextualized representations.

\medskip
\item \textbf{Engineering name:} scaled dot-product self-attention. Three questions remain --- \emph{what} gets scored, \emph{how}, and \emph{where the weights come from}. The answer to the first is $Q$, $K$, $V$.
\end{itemize}
\end{frame}

\begin{frame}{The Three Roles --- Query, Key, Value}
\begin{itemize}
\item To score and blend, each token computes \emph{three} projections of its embedding $e^{(i)} \in \mathbb{R}^{d}$:
\begin{align*}
   Q_{i} &\;=\; W_{Q}\, e^{(i)} & \text{(query)}\\
   K_{i} &\;=\; W_{K}\, e^{(i)} & \text{(key)}\\
   V_{i} &\;=\; W_{V}\, e^{(i)} & \text{(value)}
\end{align*}
where $W_{Q}, W_{K}, W_{V} \in \mathbb{R}^{k \times d}$ are learned matrices and typically $k \le d$.

\medskip
\item \textbf{The three roles} (information-retrieval mnemonic):
\begin{itemize}
    \item $Q_i$ --- the \emph{question} token $i$ asks: ``what context do I need?''
    \item $K_j$ --- the \emph{advertisement} token $j$ posts: ``here's what I'm about.''
    \item $V_j$ --- the \emph{content} token $j$ delivers if it is matched.
\end{itemize}

\medskip
\item \textbf{Why three separate projections?} Separation enables \emph{asymmetric} attention: \emph{interest} can ask a question ($Q_5$) that matches \emph{rates}'s advertisement ($K_6$) without \emph{rates} having to ask the same question back. A single shared vector could not express ``$i$ needs $j$ but not vice versa.''

\medskip
\item The match $\langle Q_i, K_j\rangle$ sets \emph{how much} to read; the value $V_j$ is \emph{what} gets read. The next slide shows this split visually.
\end{itemize}
\end{frame}

\begin{frame}{Q/K/V at a Glance --- Match Sets the Weight, the Weight Reads the Value}
\centering
\begin{tikzpicture}[
  every node/.style={font=\scriptsize},
  qbox/.style={draw, rounded corners=1.5pt, minimum width=52mm, minimum height=10mm, align=center, fill=blue!12},
  kbox/.style={draw, rounded corners=1.5pt, minimum width=40mm, minimum height=8mm, align=center, fill=orange!18},
  vbox/.style={draw, rounded corners=1.5pt, minimum width=40mm, minimum height=8mm, align=center, fill=green!18},
  outbox/.style={draw, rounded corners=1.5pt, minimum width=62mm, minimum height=9mm, align=center, fill=RedTitle!10},
  arr/.style={-{Latex[length=1.8mm]}, thick},
  score/.style={{Latex[length=1.5mm]}-{Latex[length=1.5mm]}, dashed, thick, blue!60!black}
]
\node[qbox] (Q) at (0,2.7)
    {\textbf{Query} from \emph{interest}:\ $Q_5 = W_Q e^{(5)}$\\ ``what context fixes my meaning?''};
\node[kbox] (K6) at (-3.1,0.7) {\textbf{Key} of \emph{rates}:\ $K_6$};
\node[kbox] (K9) at (3.1,0.7)  {\textbf{Key} of \emph{inflation}:\ $K_9$};
\node[vbox] (V6) at (-3.1,-0.9) {\textbf{Value} of \emph{rates}:\ $V_6$};
\node[vbox] (V9) at (3.1,-0.9)  {\textbf{Value} of \emph{inflation}:\ $V_9$};
\node[outbox] (O) at (0,-2.6)
    {blend: $0.50\,V_6 + 0.40\,V_9 + \cdots \;\to\; \tilde e^{(5)}$};
\draw[score] (Q.south west) -- node[left=1pt, font=\tiny, align=right]
    {match $\langle Q_5, K_6\rangle$\\ $\Rightarrow \alpha_6 = 0.50$} (K6.north);
\draw[score] (Q.south east) -- node[right=1pt, font=\tiny, align=left]
    {match $\langle Q_5, K_9\rangle$\\ $\Rightarrow \alpha_9 = 0.40$} (K9.north);
\draw[arr, gray!70] (K6) -- node[left, font=\tiny]{gates} (V6);
\draw[arr, gray!70] (K9) -- node[right, font=\tiny]{gates} (V9);
\draw[arr] (V6.south) -- node[left=2pt, font=\tiny]{$\times\, 0.50$} (O.north west);
\draw[arr] (V9.south) -- node[right=2pt, font=\tiny]{$\times\, 0.40$} (O.north east);
\end{tikzpicture}

\medskip
\footnotesize
\textbf{Three roles, kept separate.} The \emph{query}--\emph{key} match (dashed) only decides \emph{how much} to read; what gets read is the \emph{value}. Asymmetry matters: \emph{interest} queries \emph{rates} without \emph{rates} querying back. $W_Q, W_K, W_V$ are all learned from the loss --- nobody hand-codes ``attend to rates.''
\end{frame}

\begin{frame}{The Retrieval Analogy: A Soft Database Lookup}
\textbf{Attention is a search engine that returns a weighted blend instead of one hit.}
\begin{center}
\renewcommand{\arraystretch}{1.3}
\footnotesize
\begin{tabular}{l l l}
\toprule
\textbf{Retrieval system} & \textbf{Self-attention} & \textbf{At position $i$} \\
\midrule
search query & query vector & $Q_i = W_Q e^{(i)}$: ``what do I need?'' \\
document tags / index & key vectors & $K_j = W_K e^{(j)}$: ``what I am about'' \\
document contents & value vectors & $V_j = W_V e^{(j)}$: what is actually read \\
relevance ranking & scaled dot products & $\langle Q_i, K_j\rangle/\sqrt{k}$ \\
click the top result & softmax blend & read \emph{all} results, weighted \\
\bottomrule
\end{tabular}
\end{center}
\medskip
\begin{itemize}
\item Every position plays all three roles at once: it queries the others, advertises itself, and delivers content when matched.
\item The ``soft'' part --- a convex combination rather than a hard top-1 pick --- keeps the whole lookup differentiable, so it can be \emph{learned} end-to-end by gradient descent.
\item Preview: the same query--key logic powers document retrieval in RAG (Lec08) --- there the ``keys'' are your corpus.
\end{itemize}
\end{frame}

%=============================================================================
\section{Where Do Q, K, V Come From?}
\sectiondivider{Where Do Q, K, V Come From?}{The projection matrices are learned --- by one global objective}

\begin{frame}{Where the Projection Matrices Come From}
\small
\textbf{The $0.50/0.40$ pattern is not in the code --- it is \emph{learned}. Here is the objective that produces it.}
\begin{itemize}
\item \textbf{What is trainable ($\theta$).} In every layer $W_{Q}, W_{K}, W_{V}, W_{O}$, plus the FFN, LayerNorm, token-embedding ($\mathcal{E}$) and output-head ($W_{\text{vocab}}$) weights. All start \emph{random}, updated by gradient descent on a \textbf{single global loss}.
\smallskip
\item \textbf{The objective} (GPT-style autoregressive LM) --- predict the next token from the past:
\[
   \mathcal{L}(\theta) \;=\; -\sum_{s}\sum_{i} \log P_{\theta}\bigl(t_{s}^{(i+1)} \mid t_{s}^{(1:i)}\bigr).
\]
(BERT-style \emph{masked} LM predicts hidden tokens from both sides; fine-tuning swaps in a task label.)
\smallskip
\item \textbf{How that shapes $W_Q, W_K, W_V$.} One backward pass sends gradients through softmax, the dot products, \emph{and} the projections together. They move so that $\langle Q_i, K_j\rangle$ is \emph{high exactly where attending $i\!\to\!j$ lowers the loss}. Nobody codes ``attend to neighbours'' --- the model \emph{discovers} which pairings help predict text.
\smallskip
\item \textbf{Back to \emph{interest}.} Reading it against \emph{rates} and \emph{inflation} helped predict the surrounding words of millions of policy sentences --- \emph{that} payoff is why $W_Q, W_K, W_V$ end up producing $\alpha_6\approx0.50,\ \alpha_9\approx0.40$. Nobody wrote the rule.
\end{itemize}
\end{frame}

%=============================================================================
\section{The Mechanics, Step by Step}
\sectiondivider{The Mechanics, Step by Step}{We have Q, K, V --- now score, normalize, blend, and project back}

\begin{frame}{Mechanics I --- Scaled Dot-Product Scores}
\textbf{We already have $Q, K, V$. Now score every candidate position against \emph{interest}'s query.}
\begin{itemize}
\item Query--key compatibility for pair $(i,j)$ is the scaled dot product $\;s_{j}^{(i)} = \langle Q_{i},\, K_{j}\rangle / \sqrt{k}$.
\item \textbf{Why divide by $\sqrt{k}$?} Larger $k$ inflates dot products, pushing softmax into saturation (one term dominates, gradients vanish); dividing by $\sqrt{k}$ holds $\mathrm{Var}\bigl(s_j^{(i)}\bigr)\approx 1$.
\end{itemize}

\smallskip
\textbf{Worked example}, $i=5$ (\emph{interest}) --- the nine scaled scores:
\begin{center}\footnotesize
\renewcommand{\arraystretch}{1.15}
\begin{tabular}{l|ccccccccc}
\toprule
token $j$ & rates$_6$ & infl.$_9$ & raised$_4$ & int.$_5$ & bank$_3$ & rising$_8$ & central$_2$ & amid$_7$ & the$_1$\\
\midrule
$s_j^{(5)}$ & $2.3$ & $2.1$ & $-0.8$ & $-1.0$ & $-1.1$ & $-1.2$ & $-1.4$ & $-1.6$ & $-1.8$\\
\bottomrule
\end{tabular}
\end{center}
Only the two policy words score high; function and distant words score low.

\smallskip
\begin{itemize}
\item \textbf{Direction.} This illustration is \emph{bidirectional} (encoder-style: every token may see every other), so \emph{interest} can look ahead to \emph{rates} / \emph{inflation}. A generator would instead apply a \textbf{causal mask} --- next slide.
\end{itemize}
\end{frame}

\begin{frame}{Mechanics II --- Softmax to Attention Weights}
\begin{itemize}
\item Convert the scores into a probability distribution over positions:
\[
   \alpha_{j}^{(i)} \;=\; \frac{\exp\bigl(s_{j}^{(i)}\bigr)}{\sum_{r=1}^{n}\exp\bigl(s_{r}^{(i)}\bigr)},
   \qquad \alpha_{j}^{(i)} \ge 0, \ \ \sum_{j}\alpha_{j}^{(i)} = 1.
\]

\medskip
\item \textbf{Our example.} Exponentiate the nine scores and normalize (try it: $e^{2.3}/\sum_r e^{s_r} \approx 9.97/20.2 \approx 0.49$):
\[
   \alpha_{6}^{(5)} \approx 0.50\ (\text{\emph{rates}}), \qquad
   \alpha_{9}^{(5)} \approx 0.40\ (\text{\emph{inflation}}), \qquad
   \sum_{\text{other 7}} \alpha_{j}^{(5)} \approx 0.10.
\]
This is the $0.50/0.40$ split previewed on slide~1 --- now \emph{derived}, not asserted.
\smallskip
\item \textbf{Why softmax?} It turns any real score vector into a \emph{convex} weighting (each $\ge 0$, summing to 1) and is smooth $\Rightarrow$ differentiable end-to-end. Concentrated $\alpha \Rightarrow$ ``hard'' attention; diffuse $\alpha \Rightarrow$ ``soft.''
\smallskip
\item \textbf{Causal masking (decoder variant).} A generator masks $s_j^{(i)} = -\infty$ for $j>i$ before softmax (no peeking ahead) --- that would hide positions $6\text{--}9$ from \emph{interest}. Our bidirectional (encoder) reading keeps them.
\end{itemize}
\end{frame}

\begin{frame}{Mechanics III --- Weighted Sum of Values}
\begin{itemize}
\item Given the weights $\alpha_{j}^{(i)}$, the attention output for position $i$ is the weighted sum of \emph{value} vectors:
\[
   \text{AttnVec}_{i} \;=\; \sum_{j=1}^{n} \alpha_{j}^{(i)}\, V_{j} \;\in\; \mathbb{R}^{k}.
\]

\medskip
\item \textbf{Our example:}
\[
   \text{AttnVec}_{5} \;\approx\; 0.50\,V_{6} \;+\; 0.40\,V_{9} \;+\; 0.10\times(\text{the rest}).
\]
\emph{interest} ``reads'' mostly \emph{rates}'s and \emph{inflation}'s content. Recall the split of labour: $K_j$ decided \emph{how much}; $V_j$ is \emph{what} is delivered.

\medskip
\item \textbf{Why a blend, not a hard top-1 pick?} The output mixes $V_{\text{rates}}, V_{\text{inflation}}, \ldots$ into one vector reflecting all relevant contexts at once --- richer than any single lookup, and differentiable so it can be learned.

\medskip
\item \textbf{Contextualization, made concrete.} Swap the sentence to ``The discount \emph{rates} applied to pension funds were revised'' and the \emph{same} machinery shifts its reading: \emph{discount} and \emph{pension} become the high-weight sources, and the vector lands in discounting / pensions territory instead. Same machinery, different sentence, different meaning.
\end{itemize}
\end{frame}

\begin{frame}{Mechanics IV --- Project Back, and the Full Recipe}
\begin{itemize}
\item AttnVec lives in $\mathbb{R}^{k}$; the next layer needs $\mathbb{R}^{d}$. Project back with a learned $W_{O} \in \mathbb{R}^{d \times k}$:
\[
   \tilde{e}^{(i)} \;=\; W_{O}\,\text{AttnVec}_{i}.
\]
For \emph{interest}, $\tilde{e}^{(5)}$ is the context-aware vector we set out to build on the very first slide.

\medskip
\item \textbf{Why project back?} In multi-head attention each head works in a smaller subspace $\mathbb{R}^{k}$ ($k = d/H$); $W_O$ restores the common dimension $d$ after concatenating heads, so the result flows into the next layer with no shape change.

\medskip
\item \textbf{The full recipe (one head):}
\begin{enumerate}
    \item \textbf{Project:} $Q_i, K_i, V_i = W_Q e^{(i)},\ W_K e^{(i)},\ W_V e^{(i)}$
    \item \textbf{Score:} $s_{j}^{(i)} = \langle Q_i, K_j\rangle / \sqrt{k}$
    \item \textbf{Softmax:} $\alpha_{j}^{(i)}$
    \item \textbf{Blend:} $\text{AttnVec}_i = \sum_j \alpha_{j}^{(i)} V_j$
    \item \textbf{Project back:} $\tilde{e}^{(i)} = W_O\, \text{AttnVec}_i$
\end{enumerate}

\medskip
\item All four matrices $W_Q, W_K, W_V, W_O$ are learned by backpropagation from the task loss (previous section) --- and the whole computation runs in parallel for every position $i$ at once.
\end{itemize}
\end{frame}

\begin{frame}{Multi-Head Attention}
\small
\begin{itemize}
\item One $(Q, K, V)$ trio captures only one ``view'' of the sequence. \textbf{Multi-head attention} runs $H$ such trios in parallel and concatenates their outputs.

\medskip
\item Each head $h$ has its own projection matrices $W_{Q}^{(h)}, W_{K}^{(h)}, W_{V}^{(h)}$, all of dimension $k = d/H$.

\medskip
\item After running $H$ attention computations independently, concatenate and project:
\[
   \text{MultiHead}_{i} \;=\; W_{O}^{*}\,\Bigl[\,\text{AttnVec}_{i}^{(1)};\,\ldots;\,\text{AttnVec}_{i}^{(H)}\Bigr].
\]

\medskip
\item \textbf{Do heads have different objectives?} \emph{No --- all $H$ heads share the same training loss} (e.g.\ next-token cross-entropy). What differs is their random initialization. Specialization is \textbf{emergent}: under one global loss, $H$ parallel feature detectors diversify because identical copies of the same head would be redundant.

\medskip
\item \textbf{Observed specializations in financial text:}
\begin{itemize}
    \item Head 1: monetary-policy relations (\emph{inflation} $\leftrightarrow$ \emph{interest})
    \item Head 2: temporal phrases (\emph{Q1}, \emph{next month}, \emph{by year-end})
    \item Head 3: hawkish vs.\ dovish stance markers
\end{itemize}
Nobody told the heads what to look for --- they were not assigned separate objectives.

\medskip
\item Total parameter cost is essentially the same as a single-head attention with the full dimension $d$.
\end{itemize}
\end{frame}

\begin{frame}{Thinkbox: Attention Entropy as a Measurement Diagnostic}
\small
\textbf{The weight distribution itself carries information about measurement quality.}
\begin{itemize}
\item Position $i$'s weights $\boldsymbol\alpha^{(i)}$ form a probability distribution; its Shannon entropy $H(\boldsymbol\alpha^{(i)}) = -\sum_j \alpha_j^{(i)}\log\alpha_j^{(i)}$ measures how \emph{focused} the reading is.
\item $H \approx 0$: meaning pinned down by a few strong context words --- 2022-style tightening statements, where \emph{inflation} attends to \emph{raised}, \emph{persistent}, \emph{mandate}.
\item $H \approx \log n$: attention spread thin --- the contextual vector approaches a global average. 2019-style ``crosscurrents'' statements: \emph{inflation} attends diffusely to \emph{employment}, \emph{growth}, \emph{uncertainty}, \ldots
\end{itemize}
\begin{cautionbox}
High-entropy passages produce noisier stance scores --- measurement-error variance that \emph{moves with the state of the economy}. If you regress on Transformer-based stance measures, attention entropy is a candidate flag for heteroskedastic measurement error. A PBOC wrinkle to discuss: boilerplate political formulas share every passage with the substantive policy words --- does that push the attention entropy of terms like ``inflation'' systematically higher in PBOC reports than in FOMC statements?
\end{cautionbox}
\smallskip
\textit{\footnotesize Caveat: attention weights are intermediate computations, not structural parameters --- diagnostics and interpretation aids, never ``effects.''}
\end{frame}

%=============================================================================
\begin{frame}{Bridge: One Head Is a Building Block}
\textbf{We can now compute a context-aware vector --- once.}
\begin{itemize}
\item One attention pass mixes information \emph{across} positions. On its own it is still linear in the values, blind to word order, and single-pass.
\item \textbf{T3b wraps the head into a layer:} positional encodings (order), residual connections $+$ LayerNorm (trainability at depth), a feed-forward network (nonlinearity --- and, it turns out, the model's stored knowledge) --- then stacks the layer $L$ times and chooses among three architectural variants, ending with a full research case on 10-K filings.
\end{itemize}
\bigskip
\centering\footnotesize\textit{Arc: T1 FRAME $\to$ T2 REPRESENT $\to$ \textbf{T3 ATTEND} $\to$ T4 PRETRAIN $\to$ T5 MEASURE.}
\end{frame}

\end{document}
